\chapter{Construction du graphe technique et documentaire}

\section{Conversion du XMI}

Le format XMI contient les définitions du métamodèle mais reste difficile à inspecter et à importer directement. Le script \code{arianes/pipelines/xmi\_to\_json.py} analyse les éléments XML et produit des fichiers JSON contenant les classes, leurs attributs, leurs types et leurs relations. Cette représentation intermédiaire facilite les contrôles et découple le parseur XMI du constructeur Neo4j.

\capture{xmi-exemple.png}{Extrait du fichier XMI décrivant une classe Arianes}{fig:xmi}
\capture{json-exemple.png}{Représentation JSON obtenue après conversion}{fig:json}

\section{Modèle Neo4j}

Le constructeur crée notamment les labels \code{Classe}, \code{Attribut}, \code{Type}, \code{Enum} et \code{ValeurEnum}. Les relations structurantes sont présentées dans le tableau~\ref{tab:relations-uml}.

\begin{table}[H]
\centering
\caption{Relations principales du graphe UML}
\label{tab:relations-uml}
\begin{tabularx}{\textwidth}{p{0.28\textwidth}Y}
\toprule
Relation & Signification \\
\midrule
\fichier{A_ATTRIBUT} & Une classe possède un attribut. \\
\fichier{EST_DE_TYPE} & Un attribut possède un type. \\
\fichier{HERITE_DE} & Une classe hérite d'une autre. \\
\fichier{ASSOCIE_A} & Deux classes sont associées. \\
\fichier{COMPOSE_DE} & Composition entre classes. \\
\fichier{AGREGE} & Agrégation entre classes. \\
\fichier{A_VALEUR} & Une énumération contient une valeur. \\
\bottomrule
\end{tabularx}
\end{table}

\capture{neo4j-ligne.png}{Sous-graphe Neo4j de la classe Ligne}{fig:neo4j-ligne}

\section{Documents et chunks}

Chaque fiche devient un nœud \code{Document}. Son contenu est découpé en nœuds \code{Chunk}, reliés par \code{APPARTIENT\_A}. Le découpage produit des unités suffisamment courtes pour une recherche précise, tout en conservant l'ordre du passage dans le document.

Chaque chunk possède un embedding Mistral de dimension 1024. Un index vectoriel et un index plein texte permettent une recherche hybride. Lors de l'audit, les 244 chunks disposaient tous d'un embedding.

\section{Volumétrie}

\begin{table}[H]
\centering
\caption{État observé du graphe technique}
\begin{tabularx}{\textwidth}{Y>{\raggedleft\arraybackslash}p{1.2cm}Y>{\raggedleft\arraybackslash}p{1.2cm}}
\toprule
Élément & Nombre & Élément & Nombre \\
\midrule
Classes & 998 & Attributs & 738 \\
Documents & 33 & Chunks & 244 \\
Énumérations & 47 & Valeurs d'énumération & 146 \\
Types & 7 & Embeddings de chunks & 244 \\
\bottomrule
\end{tabularx}
\end{table}

Ces chiffres décrivent l'état testé de la base. Ils doivent être recalculés après une reconstruction du graphe.

\capture{neo4j-compteurs.png}{Résultat des requêtes de comptage et vérification des embeddings}{fig:compteurs}
